\documentclass[11pt,letter]{article}
\usepackage[T1]{fontenc}
\usepackage{lmodern}
\usepackage[margin=18mm]{geometry}
\usepackage{xcolor}
\usepackage{listings}
\usepackage{tabularx}
\usepackage{enumitem}
\usepackage{tcolorbox}
\usepackage{microtype}
\usepackage[hidelinks]{hyperref}
\hypersetup{pdftitle={How Runtime Enforcement Works},pdfauthor={}}
\definecolor{policygray}{gray}{0.93}
\setlength{\parindent}{0pt}
\setlength{\parskip}{5pt}
\setlist{nosep,leftmargin=1.5em}
\lstset{basicstyle=\ttfamily\footnotesize,columns=fullflexible,
  keepspaces=true,breaklines=true,showstringspaces=false,
  aboveskip=5pt,belowskip=5pt,xleftmargin=4pt}
\newtcolorbox{originalpolicy}{colback=policygray,colframe=policygray,
  boxrule=0pt,arc=0pt,left=7pt,right=7pt,top=5pt,bottom=5pt,
  before skip=6pt,after skip=4pt}
\newcommand{\pred}[2]{\texttt{#1}: #2\par}
\newcommand{\added}[1]{\colorbox{policygray}{%
  \makebox[\dimexpr\linewidth-2\fboxsep\relax][l]{\strut\ttfamily\small #1}}}
\begin{document}
\begin{center}
  {\LARGE\bfseries How Agent Safety Runtime Enforcement\\[2pt]
   and Inspection Work Underneath}
\end{center}
\vspace{-5pt}

We assume that the policy violations we aim to prevent occur at tool-call
boundaries, including message sends. Our tool instruments these entry points
to inspect each call, recording policy-relevant information as facts, and
enforces policy by checking each new request against rules and current
facts---for example, whether the caller may execute a transfer. Calls that
fail the checks are blocked before the tool body runs. Rules and facts are
evaluated by a Datalog engine. Datalog is a logic programming language in
which a policy is a set of rules that derive conclusions, such as whether a
request is allowed, from recorded facts.

\section{Built-in Predicate Catalog for Safety Control}


A \textbf{fact} is a recorded piece of information. For example,
\texttt{Parameter("q1", "recipient", "Bob")} records that request
\texttt{q1} has a parameter named \texttt{recipient} whose value is \texttt{Bob}.
In the general form \texttt{Parameter(Q, Name, Value)}, \texttt{Q} identifies
the request, \texttt{Name} names a parameter, and \texttt{Value} gives its value.
Our tool stores facts in a fact database (EDB).

A \textbf{predicate} names a kind of fact, such as \texttt{Parameter}.
A \textbf{rule} tells us what follows from facts we already have.
Read \texttt{:-} as ``if'' and commas on its right as ``and'': the conclusion
on the left holds when all conditions on the right hold. Letters such as
\texttt{P} and \texttt{R} stand for values; the same letter within a rule
refers to the same value.

\subsection*{Runtime control for actions}
Every action an agent takes, such as a tool call or a message send, passes
through a gate that decides whether it may run; because the check happens at
the moment of action, it applies no matter how the agent was prompted or what
content it has seen.

\pred{Request(Q, Principal, EventKind)}{Request \texttt{Q} is made by
\texttt{Principal} to perform \texttt{EventKind}.}
\pred{Parameter(Q, Name, Value)}{Request \texttt{Q} has this parameter and value.}
\pred{Allow(Q)}{The rules permit request \texttt{Q}.}
\pred{Deny(Q, Reason)}{The rules reject request \texttt{Q} for this reason.}
\begin{center}
\begin{tabular}{lll}
\textbf{Allow(Q)} & \textbf{Any Deny(Q, Reason)} & \textbf{Tool body} \\
\hline
Absent & Absent & Blocked \\
Present & Absent & Allowed to run \\
Absent & Present & Blocked \\
Present & Present & Blocked
\end{tabular}
\end{center}
Our tool checks each request using the current facts. A request needs
\texttt{Allow}; any \texttt{Deny} blocks it.


\subsection*{Category 1: Capability and access control}
Access control dates back to early operating-system protection. The
group-based rules below follow later work on role-based access control and
Datalog authorization languages such as Binder (2002); here they decide which
agent may invoke which tool on which resource.

\pred{Capability(Principal, Scope, EventKind)}{This principal has permission
to perform this operation on this resource or group of resources.}
\pred{PrincipalGroup(Principal, Group)}{This principal belongs to this group.}
\pred{ResourceGroup(Parent, Child)}{This resource belongs to a larger resource group.}
\pred{Can(Principal, Resource, EventKind)}{This principal has the permission,
directly or through a group.}

\begin{lstlisting}
Can(P, R, E) :- Capability(P, R, E).
Can(P, R, E) :- PrincipalGroup(P, G), Can(G, R, E).
Can(P, Child, E) :- ResourceGroup(Parent, Child), Can(P, Parent, E).
\end{lstlisting}
Read each rule as an if-then sentence:
\begin{enumerate}
\item If \texttt{P} has the capability, \texttt{P} can perform \texttt{E} on \texttt{R}.
\item If \texttt{P} is in group \texttt{G} and \texttt{G} can, so can \texttt{P}.
\item If \texttt{P} can act on \texttt{Parent}, it can act on each \texttt{Child} in it.
\end{enumerate}
Rules 2 and 3 are \emph{recursive}: \texttt{Can} appears on both sides, so
permissions pass through groups nested to any depth. The engine reapplies the
rules until no new facts appear, which always terminates. The app's
\texttt{Allow} rules use \texttt{Can} to check the caller's permission.

\subsection*{Category 2: Information flow control}
Information flow control dates back to Denning's lattice model (1976) and
later taint tracking. For agents, it tracks whether untrusted content, such
as an external message, influenced a later action.

\pred{Source(Principal, Data)}{This data came directly from this principal.}
\pred{DependsOn(Output, Input)}{This output was produced using this input.}
\pred{Provenance(Data, Principal)}{This data came from this principal,
directly or through other data.}

\begin{lstlisting}
Provenance(D, P) :- Source(P, D).
Provenance(Out, P) :- DependsOn(Out, In), Provenance(In, P).
\end{lstlisting}
Read each rule as an if-then sentence:
\begin{enumerate}
\item If \texttt{D} came directly from \texttt{P}, its provenance includes \texttt{P}.
\item If \texttt{Out} was produced from \texttt{In}, it inherits every origin of \texttt{In}.
\end{enumerate}
Rule 2 is recursive, so origins carry forward through any chain of copying,
combining or rewriting. The app records \texttt{Source} and \texttt{DependsOn};
the engine derives the rest.

\begin{originalpolicy}
\textbf{Example 1. Policy:} Only Payments may execute transfers. Payment details from
external messages require additional verification.
\end{originalpolicy}
The transfer rule checks the caller's \texttt{Can} permission and uses
\texttt{Provenance} to require verification for external-message details.
As long as the app records \texttt{DependsOn} for derived data, rewriting
those details keeps their origin and cannot bypass the check.


\subsection*{Category 3: Ordering and lifecycle constraints with stateful effects}
Ordering constraints build on state machines and typestate (Strom and Yemini,
1986). For agents, they ensure that steps happen in the required order, such
as approval before publication, even when different agents perform them.

\pred{ResourceType(Resource, Type)}{What kind of resource this is.}
\pred{ResourceState(Resource, State)}{What state the resource is in now.}
\pred{InitialState(Type, State)}{The starting state for a new resource of this type.}
\pred{AllowedTransition(Type, EventKind, From, To)}{This operation may move
this type of resource from one state to another.}
\pred{ReachableState(Type, From, To)}{One or more allowed steps connect these states.}
\pred{CanTransition(Resource, EventKind, From, To)}{This resource can take
this step from its current state.}

\begin{lstlisting}
ReachableState(T, From, To) :- AllowedTransition(T, _, From, To).
ReachableState(T, From, To) :- ReachableState(T, From, Mid),
                             AllowedTransition(T, _, Mid, To).
CanTransition(R, E, From, To) :- ResourceType(R, T),
    ResourceState(R, From), AllowedTransition(T, E, From, To).
\end{lstlisting}
Read each rule as an if-then sentence:
\begin{enumerate}
\item If one allowed step leads from \texttt{From} to \texttt{To}, \texttt{To} is reachable.
\item If \texttt{Mid} is reachable and one step leads from \texttt{Mid} to \texttt{To}, so is \texttt{To}.
\item If \texttt{R} is now in \texttt{From} and its type allows \texttt{E} from
\texttt{From} to \texttt{To}, \texttt{R} can take that step.
\end{enumerate}
Rule 2 is recursive, so \texttt{ReachableState} connects states through any
number of steps. The app's \texttt{Allow} rules use \texttt{CanTransition} to
check the next step. This is the stateful effect: a new resource starts in
its \texttt{InitialState}, and when an operation succeeds, our tool
\emph{replaces} its \texttt{ResourceState} with the new state rather than
adding a second one.

\begin{originalpolicy}
\textbf{Example 2. Policy:} A document may be published only after approval.
\end{originalpolicy}
Configure \texttt{DRAFT} $\rightarrow$ \texttt{APPROVED} $\rightarrow$
\texttt{PUBLISHED}. A publish rule requires the
\texttt{APPROVED} $\rightarrow$ \texttt{PUBLISHED} transition, so publishing a
draft is blocked, and a successful publication moves the document to
\texttt{PUBLISHED}.

\section{Translation: Injecting Enforcement into Your App}

Enforcement is added to an existing app through two mechanisms: code
injection and call interception. Injection inserts small additions into the
app's source, namely a \texttt{@policy} annotation on each tool function and
fact-writing statements inside it, while the policy itself stays in a separate
\texttt{.dl} file. At run time, the annotation intercepts every call before the
function body executes. It records the call as \texttt{Request} and
\texttt{Parameter} facts, asks the Datalog engine whether \texttt{Allow} holds
and no \texttt{Deny} holds for that request, and then either lets the body run
or blocks the call. When the body completes, the injected fact writes add its
results to the EDB, so the engine can evaluate later requests against what has
already happened. Because the rules live in Datalog rather than in the tool
code, you can change a policy without editing the functions it governs.

More specifically, our tool adds three pieces:
\begin{enumerate}
\item \textbf{A \texttt{.dl} file} containing application rules and base facts,
such as who has permission and which state changes are allowed.
\item \textbf{\texttt{@policy} on tool functions.} Before an annotated function
runs, our tool records \texttt{Request} and \texttt{Parameter}, checks the
rules, and blocks calls that are not permitted.
\item \textbf{Fact writes inside tool functions.} Added statements record known
application information and results in the EDB for later checks.
\end{enumerate}

Grey lines show the additions to a message-send function:

\begingroup
\setlength{\tabcolsep}{0pt}
\begin{tabularx}{\linewidth}{@{}X@{}}
\added{@policy} \\
\texttt{\small def send\_message(recipient, body):} \\
\texttt{\small\hspace*{1.5em}result = transport.send(recipient, body)} \\
\added{\hspace*{1.5em}policy.put("Sent", result.id, recipient)} \\
\texttt{\small\hspace*{1.5em}return result}
\end{tabularx}
\endgroup

\texttt{@policy} checks the caller and recipient before sending.
After a successful send, \texttt{policy.put} records which message was sent to
which recipient. Later requests can be checked against that fact.

As an end user of this framework, you do not need to understand these
internals to use it; we include the technical details so that you can see
what is happening underneath when a call is checked, allowed or blocked.

\end{document}